% F085｜真实数学图；代表函数不构成抽象解的唯一性证明。
\begin{center}
\begin{tikzpicture}[x=2.6cm,y=2.6cm,>=stealth]
\draw[->,black!65] (-.3,0)--(1.5,0) node[right,font=\small] {横轴};
\draw[->,black!65] (0,-.15)--(0,1.3);
\draw[CExplanation,thick] (0,0)--(30:1) node[right,font=\small,xshift=4pt] {$P_-$};
\draw[CExplanation,thick] (0,0)--(60:1) node[above right,font=\small] {$M$};
\draw[CExplanation,thick] (0,0)--(90:1) node[above right,font=\small] {$P_+$};
\draw[CExplanation,thick] (0:1) arc (0:90:1);
\draw[orange!80!black,very thick] (30:1)--(90:1);
\fill[orange!80!black] (.4330127,.75) circle (.025);
\draw[orange!80!black] (.47,.78)--(1.1,1.15) node[right,font=\small] {$Q$（弦中点）};
\draw[dashed] (30:1)--(.8660254,-.25);
\draw[dashed] (90:1)--(0,-.25);
\draw[dashed] (.4330127,.75)--(.4330127,-.25);
\draw[<->,thick,orange!80!black] (0,-.18)--(.8660254,-.18);
\node[below,font=\small] at (.433,-.25) {两端横坐标的平均值是 $Q$ 的横坐标};
\draw (0.23,0) arc (0:60:.23);
\node[font=\small] at (.25,.26) {$x$};
\node[font=\small,below left] at (0,0) {$O$};
\node[font=\small,align=left,anchor=west] at (1.9,.58) {$P_\pm=(\cos(x\pm y),\sin(x\pm y))$\\[5pt]
$OM=1,\quad OQ=\cos y$\\[5pt]
$Q$ 在方向角 $x$ 的射线上\\[5pt]
$Q$ 的横坐标 $=\cos x\cos y$};
\end{tikzpicture}
\par\small 图取 $x=\pi/3,y=\pi/6$。单位圆的弦中点给出 $[\cos(x+y)+\cos(x-y)]/2=\cos x\cos y$。
\end{center}
